\subsection{Topological Spaces, Bases, and Separation}

% **Role:** abstract structure behind limits and continuity.
%
% Topological spaces, bases, product topology, quotient topology, countability axioms.


General topology abstracts the open-set behavior of metric spaces. This selective pass records spaces, bases, products, quotients, and countability axioms for later use.

\begin{definition} \label{def:topological-space} \hlt{Topological Space}\\
A \hlt{topology} on a set $X$ is a collection $\mathcal{T}\subseteq 2^X$ such that:
\begin{enumerate}[label=(\roman*)]
\setlength{\itemsep}{0pt}
\item $\emptyset,X\in\mathcal{T}$;
\item if $\{\mathcal{O}_\alpha\}_{\alpha\in A}\subseteq\mathcal{T}$, then $\bigcup_{\alpha\in A}\mathcal{O}_\alpha\in\mathcal{T}$;
\item if $\mathcal{O}_1,\ldots,\mathcal{O}_n\in\mathcal{T}$ with $n\geq1$, then $\bigcap_{j=1}^n\mathcal{O}_j\in\mathcal{T}$.
\end{enumerate}
Then $(X,\mathcal{T})$ is a \hlt{topological space}. A set $U\subseteq X$ is \hlt{open} iff $U\in\mathcal{T}$, and $F\subseteq X$ is \hlt{closed} iff $X\setminus F\in\mathcal{T}$.
\end{definition}

\begin{proposition} \label{prop:metric-open-sets-form-topology} \hlt{Metric Open Sets Form a Topology}\\
Let $(X,d)$ be metric space, $B_r(x)\defeq\{y\in X:d(x,y)<r\}$, and $\tau_d\defeq\{U\subseteq X:\forall x\in U,\ \exists r>0,\ B_r(x)\subseteq U\}$. Then $\tau_d$ is a topology on $X$.
\begin{proof}
Clearly $\emptyset,X\in\tau_d$. If $\{\mathcal{O}_\alpha\}_{\alpha\in A}\subseteq\tau_d$ and $x\in\bigcup_{\alpha\in A}\mathcal{O}_\alpha$, choose $\alpha$ with $x\in\mathcal{O}_\alpha$ and $r>0$ with $B_r(x)\subseteq\mathcal{O}_\alpha$; hence $B_r(x)\subseteq\bigcup_{\alpha\in A}\mathcal{O}_\alpha$. If $\mathcal{O}_1,\ldots,\mathcal{O}_n\in\tau_d$, $n\geq1$, and $x\in\bigcap_{j=1}^n\mathcal{O}_j$, choose $r_j>0$ with $B_{r_j}(x)\subseteq\mathcal{O}_j$ and set $r\defeq\min_{1\leq j\leq n} r_j$; then $B_r(x)\subseteq\bigcap_{j=1}^n\mathcal{O}_j$. Thus $\tau_d$ is a topology on $X$.\\
\hlt{Structural Commentary}: Metric balls generate an open-set system with the exact axioms of a topology.\\
\hlt{Key Insight}: General topology keeps the open-set calculus and forgets the distance function.
\end{proof}
\end{proposition}

\begin{definition} \label{def:basis-for-topology} \hlt{Basis}\\
Let $X$ be a set. A collection $\mathcal{B}\subseteq 2^X$ is a \hlt{basis} for a topology on $X$ iff:
\begin{enumerate}[label=(\roman*)]
\setlength{\itemsep}{0pt}
\item $X=\bigcup_{B\in\mathcal{B}}B$;
\item if $B_1,B_2\in\mathcal{B}$ and $x\in B_1\cap B_2$, then there exists $B_3\in\mathcal{B}$ with $x\in B_3\subseteq B_1\cap B_2$.
\end{enumerate}
The \hlt{topology generated by $\mathcal{B}$} is $\mathcal{T}_{\mathcal{B}}\defeq\{\bigcup_{B\in\mathcal{U}}B:\mathcal{U}\subseteq\mathcal{B}\}$.\\
A collection $\mathcal{S}\subseteq 2^X$ is a \hlt{subbasis} iff $\{\bigcap_{j=1}^n S_j:n\geq1,\ S_1,\ldots,S_n\in\mathcal{S}\}$ is a basis.
\end{definition}

\begin{example} \label{ex:metric-balls-basis} \hlt{Metric Balls as a Basis}\\
In a metric space, the collection of all open balls is a basis for the metric topology. On $\R^k$, balls with rational centers and rational radii form a countable basis.
\end{example}
